\section{中文编译（三）：中国 AI 产业与西方实验室的差异}

\subsection{AI 模型公司已经不只是研究团队}

\begin{translationbox}{原文主线编译}
作者说，今天构建 AI 模型最有趣的地方在于，它不再只是把一群优秀研究者聚在一栋楼里，做出一个工程奇迹。过去也许可以这样理解，但要维持 AI business，LLM 正变成构建、部署、融资和采用的混合体。领先 AI 公司存在于复杂生态中，需要资金、算力、数据等持续输入，才能继续推动前沿。
\end{translationbox}

这段把讨论从``研究者心态''转向``产业系统''。OpenAI、Anthropic 这样的公司已经是资本、云基础设施、数据供应链、企业市场、消费产品和政策讨论的交汇点。中国实验室也面对相同问题，但它们的解法不同。

\subsection{国内 AI 需求的早期信号}

\begin{translationbox}{原文主线编译}
作者提到一个常见假设：中国 AI 市场会更小，因为中国公司通常不愿意为软件付费，因此无法形成支撑实验室的大规模推理市场。作者认为，这个说法只适用于映射到 SaaS 生态的软件支出；而中国同时显然有很大的云市场。关键问题是：企业 AI 支出到底会像 SaaS 一样很小，还是像 cloud 一样成为基础设施？中国实验室内部也在争论这个问题。作者的总体感受是，AI 更像 cloud，而且没有人特别担心围绕新工具的市场不会成长。
\end{translationbox}

\begin{importantbox}{SaaS 逻辑 vs. Cloud 逻辑}
如果 AI 是 SaaS，中国市场可能受制于企业软件付费传统；如果 AI 是 cloud，它就是生产力基础设施，企业会为算力、推理、自动化和业务能力付费。作者倾向于认为 AI 更接近后者。
\end{importantbox}

这点对判断中国模型公司收入很重要。很多外部分析会说：中国企业不买软件，所以中国大模型 API 赚不到钱。但如果 AI 能力被嵌入云、办公、客服、研发、推荐、营销、设备和本地生活系统，它未必走传统 SaaS 路径。

\subsection{中国开发者深受 Claude 影响}

\begin{translationbox}{原文主线编译}
作者发现，中国许多 AI 开发者都非常迷 Claude，认为 Claude 改变了他们构建软件的方式，尽管 Claude 名义上在中国不可用。作者说，中国过去不太愿意购买软件这一事实，并没有让他觉得中国不会出现大规模推理需求。中国技术人员非常务实、谦逊、有动力，这种力量似乎强于过去不买软件的习惯。也有研究者提到用自己的工具，例如 Kimi 或 GLM CLI，但几乎所有人都会提到 Claude。相比之下，Codex 被提到的次数少得惊人，而 Codex 在 Bay Area 正明显变得流行。
\end{translationbox}

这段有两个信息：第一，中国开发者会绕过可访问性障碍使用最有效工具；第二，coding agent 的真实需求可能比官方市场数据更早出现。

\begin{knowledgebox}{为什么 Claude 是需求信号}
开发者自发使用 Claude，说明 coding assistant / coding agent 已经从玩具变成生产工具。即使官方渠道受限，强需求也会通过替代路径出现。对模型公司来说，这意味着未来推理需求不只来自聊天机器人，而来自软件开发工作流本身。
\end{knowledgebox}

\subsection{中国公司的技术所有权心态}

\begin{translationbox}{原文主线编译}
作者说，中国文化和强劲经济引擎结合，产生了很难预测的结果。他留下的持久印象是，众多 AI 模型并不是某个总体规划的结果，而是许多科技公司在当前均衡下的务实选择。行业尊重字节和阿里这样的 incumbent，认为它们拥有资源，可能赢下很多市场。DeepSeek 是受尊重的技术领导者，但远不是市场领导者；它设定方向，却不一定具备经济上最终取胜的结构。
\end{translationbox}

这段是全文产业判断的中心：\textbf{中国模型多，不是因为所有公司都盲目追热点，而是因为很多公司认为必须掌握这项底层技术。}

\begin{translationbox}{原文主线编译}
这解释了为什么像美团或蚂蚁这样的公司也在构建模型。西方观察者可能会惊讶：为什么这些公司要做 LLM？但在中国语境下，它们显然认为 LLM 将成为未来技术产品的核心，因此需要一个强底座。当它们 fine-tune 强通用模型时，开放社区的反馈会让它们的 stack 更坚固，同时它们也可以保留内部 fine-tuned 版本用于自己的产品。行业里的 open-first mentality 很大程度上由实用主义定义：它帮助模型获得反馈，回馈开源社区，也增强公司使命。
\end{translationbox}

\begin{importantbox}{open-first 的实用主义}
开放权重不一定来自纯粹开源理想主义。对很多中国公司而言，开放是一个产品和研发策略：让社区测试模型、扩大开发者生态、建立技术声誉，同时把真正业务相关的 fine-tuned 版本留在内部。
\end{importantbox}

\subsection{政府帮助真实存在，但难以量化}

\begin{translationbox}{原文主线编译}
作者说，外界经常断言中国政府正在积极帮助开放 LLM 竞赛。政府确实是多层级、分散的，每一层并没有清晰剧本。北京不同街区会竞争吸引科技公司办公。政府提供的``帮助''几乎肯定包括减少许可证等官僚摩擦，但还能走多远？能否帮助吸引人才？能否帮助走私芯片？访问中，作者听到许多关于政府兴趣或帮助的提法，但细节太少，不足以形成确定叙事。至少，他没有看到中国最高层影响任何模型技术决策的迹象。
\end{translationbox}

这段很重要，因为它同时反对两个极端：一是认为中国 AI 全靠政府指挥；二是认为政府因素完全不存在。作者看到的是一种更复杂的地方政府、产业园区、政策兴趣和资源协调网络。

\begin{warningbox}{不要把政府想象成单一 actor}
中国政府支持 AI，可能表现为地方招商、审批便利、算力资源协调、人才政策和产业基金，但这不等于中央直接决定模型架构、RL recipe 或开源策略。产业政策影响路径往往比外部叙事复杂。
\end{warningbox}

\subsection{数据产业不如美国成熟}

\begin{translationbox}{原文主线编译}
作者此前听说 Anthropic 或 OpenAI 可能为单个环境支付 1000 万美元以上，每年为推动 RL 前沿的数据和环境花费数亿美元。因此他很想知道，中国实验室是否购买同样的美国环境，或者是否有镜像的国内数据生态。答案不是完全没有数据产业，而是中国实验室的经验是：数据产业质量相对较差，很多时候自己构建环境或数据更好。研究者自己会花大量时间制作 RL training environments，一些大公司如字节和阿里有内部数据标注团队。这与前面提到的 build-not-buy 心态一致。
\end{translationbox}

\begin{knowledgebox}{RL environment 是什么}
在 agentic RL 中，环境不是一堆静态文本，而是模型可以交互、尝试、失败、得到反馈的任务世界。例如 coding benchmark、网页操作、工具调用、数学证明、企业工作流模拟等。高质量环境越接近真实任务，越能训练出可用 agent。
\end{knowledgebox}

这里有一个有趣反转：美国有更成熟的外部供应链，中国供应链弱，因此中国团队被迫自建。短期看这是成本，长期看可能形成内部能力。

\subsection{对 Nvidia 芯片的渴望}

\begin{translationbox}{原文主线编译}
作者说，Nvidia 计算仍然是训练黄金标准，每个人的进展都受限于没有更多 Nvidia 芯片。如果供应存在，中国实验室显然会购买。其他加速器，包括但不限于华为，在推理方面被积极评价；许多实验室都能接触华为芯片。
\end{translationbox}

\begin{warningbox}{推理替代不等于训练替代}
国产或替代加速器可以在 inference 中发挥越来越大作用，但前沿训练对集群稳定性、通信、kernel、调试、框架生态和工程经验要求极高。Nvidia 仍然是训练前沿模型的事实标准。
\end{warningbox}

\subsection{中国生态是否会产生不同模型}

\begin{translationbox}{原文主线编译}
作者总结说，这些点描绘了一个非常不同的 AI 生态。如果把西方实验室的运作方式快速映射到中国同行身上，常常会犯类别错误。关键问题是：这些不同生态会不会生产出有意义地不同的模型？还是说，中国模型永远会被解释为类似美国前沿模型 3--9 个月前的版本？
\end{translationbox}

\subsection{本章小结}

第二大部分讲的是产业：AI 公司需要资金、算力、数据、部署和采用；中国公司的差异在于更强控栈、更务实开放、外部数据供应链更弱、Nvidia 训练算力更紧缺，以及可能被低估的国内推理需求。


